\item As the case shows, legal protection of publicity rights arises both from
state statutes and from the common law. When this happens, understand that they
are two separate causes of action and theories of liability. Why have both? The
exact contours are not identical, as the case illustrates. Most importantly
(though not of relevance to this case), the common law right of publicity
typically ends upon death, while state statutes often apply for years or decades
postmortem.

\item Because the case is being decided on a motion to dismiss, the court does
not identify ``facts corresponding to each legal element'' of the right of
publicity. But you should. Imagine you are counsel for the plaintiffs---what
facts would you hope to find to prove each of the elements of the legal rules?
What would you ask for during discovery to locate these facts?

\item The court rejects the argument that Clearview does not violate the
California statutory right of publicity because Clearview does not use its
database ``for purposes of advertising, selling, or soliciting.'' In rebutting
this argument, the court highlights the phrase ``in any manner'' in the statute.
Read the statutory text \emph{extremely} carefully (consider diagramming the
sentence). There's a good argument that the court highlighted the wrong phrase.
But there's also a good, different response to Clearview's argument.

\item What First Amendment or
freedom of expression concerns can you see with the right of publicity?
Consider, in particular, how publicity rights might be asserted against other
potential defendants such as journalists. What elements of the common law or
statutory tests might help to alleviate any First Amendment conflicts? How do
you think courts should interpret those elements?

\defbook{zirpoli-ai-rop}{
author=Christopher T. Zirpoli,
instauth=Congressional Research Service,
title={Artificial Intelligence Prompts Renewed
Consideration of a Federal Right of Publicity},
year=jan 29 2024,
serial=Report No. LSB11052,
url=https://www.congress.gov/crs-product/LSB11052,
}

\having{white-v-samsung}{}{}{
\defcase{white-v-samsung}{
White v. Samsung Electronics America, Inc.,
971 F.2d 1395 (9th Cir. 1992)
}
\item In \inline{white-v-samsung}\optclause, the Ninth Circuit considered
whether the publicity rights of Vanna White, the famed letter-turner of the
television game show \emph{Wheel of Fortune}, had been violated by an
advertisement that featured a robotic letter-turner in the same game show. The
humanoid robot plainly referenced White (it wore a wig, gown, and
jewelry ``selected to resemble White's hair and dress''), but the
advertisement did not identify Vanna White by name and the robot did not
resemble White's face.
\endgraf
Applying the same common-law test as \emph{Clearview}, the Ninth Circuit agreed
that ``the robot ad did not make use of White's name or likeness.'' But that did
not stop the court from finding a plausible violation:
\begin{quote}
The right of publicity does not require that appropriations of identity be
accomplished through particular means to be actionable.\ldots Indeed, if we
treated the means of appropriation as dispositive in our analysis of the right
of publicity, we would not only weaken the right but effectively eviscerate it.
The right would fail to protect those plaintiffs most in need of its protection.
Advertisers use celebrities to promote their products. The more popular the
celebrity, the greater the number of people who recognize her, and the greater
the visibility for the product. The identities of the most popular celebrities
are not only the most attractive for advertisers, but also the easiest to evoke
without resorting to obvious means such as name, likeness, or voice.
\end{quote}
\sentence{white-v-samsung at 1398-1399}. Should the right of publicity reach
this far, to anything that ``evoke[s]'' a celebrity's identity? Have you ever
evoked a celebrity's identity when conversing with friends? Which elements of
the rule, if any, avoid you and the everyone else in the world from being
liable?
}

\item Traditionally, the right of publicity was a protection for celebrities,
preventing unscrupulous advertisers from misappropriating their fame for product
advertisements and such. But like other legal devices, publicity rights are a
legal tool that may be adapted to other uses. Such rights are at the core of
name, image, and likeness (NIL) rules for student athletes. And in a world of
social media and generative artificial intelligence, lawmakers have looked to
the right of publicity as a possible response to deepfake images and privacy
invasions. \sentence{see zirpoli-ai-rop}.

What do you think of the use of old laws, intended for a different purpose,
being used to solve new problems? On the one hand, misalignments can easily crop
up, rendering legal rules unbalanced or overly focused on irrelevant matters.
But on the other hand, how else should law keep up with a rapidly changing
social and technological environment?

